\section{Kiểm tra tính cân bằng}

Trước khi chia dataset, thực hiện kiểm tra \textbf{tính cân bằng} của dữ liệu. Dataset cuối cùng sau các bước tiền xử lý có \textit{68,604 mẫu} với phân phối nhãn:

\begin{itemize}
    \item \textbf{Tin thật (label = 1):} 34,526 mẫu (50.3\%)
    \item \textbf{Tin giả (label = 0):} 34,078 mẫu (49.7\%)
\end{itemize}

\noindent Dataset cho thấy \textbf{tính cân bằng rất tốt} với sự chênh lệch chỉ 0.6\% giữa hai lớp. Điều này là một \textit{lợi thế quan trọng} vì:
\begin{itemize}
    \item Tránh được vấn đề \textit{class imbalance} trong machine learning
    \item Model sẽ không thiên lệch về một lớp cụ thể
    \item Các metric đánh giá (accuracy, precision, recall) sẽ đáng tin cậy
    \item Không cần áp dụng các kỹ thuật resampling phức tạp
\end{itemize}

\section{Chia tập train và test}

Việc chia dataset được thực hiện \textbf{trước bất kỳ bước tiền xử lý nào} để đảm bảo tính toàn vẹn của quá trình đánh giá. Chúng tôi sử dụng \texttt{train\_test\_split} với các tham số:

\begin{itemize}
    \item \textbf{Test size:} 20\% (13,721 mẫu test, 54,883 mẫu train)
    \item \textbf{Stratified sampling:} Đảm bảo tỷ lệ tin thật/giả được duy trì trong cả train và test set
\end{itemize}

\noindent Việc chia dataset \textit{ngay từ đầu, trước mọi bước preprocessing} là \textbf{cực kỳ quan trọng} để tránh \textit{data leakage}:

\begin{itemize}
    \item \textbf{Ngăn chặn information leakage:} Thông tin từ test set không được "rò rỉ" vào quá trình training qua các bước feature engineering, scaling, hoặc vectorization
    \item \textbf{Đảm bảo đánh giá chính xác:} Test set giữ nguyên trạng thái "unseen" để đánh giá hiệu suất thực tế của model
    \item \textbf{Mô phỏng production environment:} Trong thực tế, model sẽ gặp dữ liệu hoàn toàn mới chưa được xử lý
    \item \textbf{Tránh overfitting ẩn:} Các bước preprocessing như TF-IDF fitting, feature scaling được thực hiện độc lập cho train và test
\end{itemize}

\noindent Kết quả chia cho thấy cả train set (54,883 mẫu) và test set (13,721 mẫu) đều \textbf{duy trì tỷ lệ cân bằng} giữa tin thật và tin giả, đảm bảo tính đại diện và khả năng generalization của model.

\noindent Từ bước này trở đi, mọi quá trình tiền xử lý, feature engineering, và training sẽ được thực hiện \textit{riêng biệt} cho train và test set để đảm bảo tính khách quan trong đánh giá.